\input{template.tex}

\ifPDFTeX
  \renewcommand{\rmdefault}{cmr}
\fi

\def\SheetNo{4}
\setlength{\emergencystretch}{2em}

\begin{document}
\sheettitle{Логика и комбинаторика.}

\begin{problems}

\item Население острова состоит из $100$ рыцарей и $100$ лжецов.
Рыцари всегда правдивы, лжецы всегда лгут. Дружба взаимна, и у каждого
жителя есть хотя бы один друг. При опросе $100$ жителей утверждали,
что среди их друзей нет лжецов, а остальные $100$ --- что среди их
друзей нет рыцарей. Каково наименьшее возможное число дружеских пар,
в которых один житель --- рыцарь, а другой --- лжец?

\item В волшебном городе есть неограниченное количество баров. Есть $5$ практиков, у которых есть $5$ предпочитаемых баров. У каждого из $m$ студентов тоже есть $5$ любимых баров в этом городе. При каком наибольшем $m$ каждый (студенты и практики) может пойти в один из своих любимых баров так, чтобы ни в одном баре не встретились практик и студент, вне зависимости от предпочтений? Требуется показать, что при большем $m$ можно подобрать предпочтения, при которых условие невыполнимо, а при найденном $m$ предложить способ для любых предпочтений.

\item Являются ли полными следующие системы булевых функций?
\begin{parts}
  \item $\{\neg,\vee\}$;
  \item $\{\oplus,\vee\}$;
  \item $\{\oplus,\supset\}$.
\end{parts}

\item Докажите тождество
\[
  C_n^0+C_{n+1}^1+C_{n+2}^2+\dots+C_{n+k}^k=C_{n+k+1}^k.
\]

\item Сто мудрецов ставят в ряд и каждому надевают либо красный,
либо синий колпак; колпаков каждого цвета имеется по $100$.
Мудрец видит только колпаки стоящих перед ним. Затем мудрецы по одному,
в выбранном ими порядке, объявляют предполагаемый цвет собственного
колпака. Все слышат предыдущие ответы; других сигналов подавать нельзя.
Тот, кто ошибся, будет казнён. До начала испытания мудрецы могут
согласовать стратегию. Какое максимальное число спасённых они могут
обеспечить при любом распределении колпаков?

\item Между некоторыми жителями острова установлены взаимные дружеские
отношения. Каждый житель --- либо неизменно правдивый рыцарь, либо
неизменно лгущий лжец. Все жители утверждают, что рыцари составляют
больше половины их друзей. Докажите, что среди всех дружеских пар
не меньше половины составляют пары из двух рыцарей или двух лжецов.

\item Курс лекций состоял из $35$ пар. На каждой паре присутствовало ровно $27$ студентов. Оказалось, что на любых трёх занятиях был ровно один общий студент. Верно ли, что был студент, посетивший все пары?

\item Докажите формулу (а) алгебраически и (б) комбинаторно:
\[
  C_n^k C_k^l = C_n^l C_{n-l}^{k-l}.
\]

\item Среди трёх подозреваемых --- Брауна, Джонса и Смита --- ровно один
совершил преступление. Каждый сообщил следователю два утверждения.
Браун отрицает как свою причастность, так и причастность Джонса.
Смит отрицает свою причастность и обвиняет Брауна.
Джонс отрицает причастность Брауна и обвиняет Смита.
Известно, что у одного подозреваемого оба утверждения истинны,
у другого оба ложны, а у третьего одно истинно и одно ложно.
Определите виновного.

\item На острове рыцари никогда не лгут, а лжецы никогда не говорят правды.
В хороводе из $21$ жителя представлены оба типа островитян.
Каждый участник утверждает, что его два соседа принадлежат к разным типам:
один из них рыцарь, другой --- лжец. Определите число рыцарей в хороводе.

\item Не пользуясь таблицами истинности, приведите к СКНФ и СДНФ следующие функции:
\begin{parts}
  \item $((P\supset Q)\supset(R\supset\neg P))\supset(\neg Q\supset\neg R)$;
  \item $(((((P\supset Q)\supset\neg P)\supset\neg Q)\supset\neg R)\supset R)$.
\end{parts}

\item На острове $2023$ жителя, каждый из которых --- либо рыцарь,
всегда отвечающий правдиво, либо лжец, всегда дающий ложный ответ.
Антрополог спросил каждого жителя о числе его друзей на острове;
дружба взаимна. Полученные ответы --- целые числа от $0$ до $2022$,
причём одинаковые ответы допускаются. По результатам опроса антрополог
смог доказать, что число лжецов не меньше $k$. Найдите наибольшее $k$,
для которого такая ситуация возможна.

\item
\begin{enumerate}
  \item[а)] В ряд выстроены $n$ школьников. Сколькими способами можно выбрать из них $k$ дежурных так, чтобы никакие два дежурных не стояли рядом?
  \item[б)] За круглым столом короля Артура сидят $n$ рыцарей. Каждый из них враждует со своими соседями. Король хочет составить отряд из $k$ рыцарей. Сколькими способами это можно сделать так, чтобы в этом отряде не было врагов?
\end{enumerate}

\item Докажите, что число $\displaystyle\frac{(n^2)!}{(n!)^{n+1}}$ является целым для натуральных значений $n$.

\item Найдите количество способов, которыми два человека могут разделить $2n$ предметов одного сорта, $2n$ предметов второго сорта и $2n$ предметов третьего сорта так, чтобы каждый получил $3n$ предметов.

\item В Министерстве Пропаганды работали честные сотрудники (которые всегда говорили правду) и патологические лжецы (которые всегда лгали). Каждый сотрудник знал про любого другого, честный тот или лжец. Вновь назначенный Министр Пропаганды у каждого подчинённого спросил про какого-то другого подчинённого, честный тот или лжец. При этом ни про какого подчинённого Министр не спрашивал дважды. После опроса Министр выгнал из Министерства всех сотрудников, которые были названы честными. Могло ли у Министра остаться ровно $2013$ подчинённых?

\item
\begin{enumerate}
  \item[а)] (Формула Вандермонда.) Пусть $m,n,k\in\mathbb N$, причём $m,n\geq k$. Докажите, что
  \[
    C_n^0C_m^k+C_n^1C_m^{k-1}+\dots+C_n^kC_m^0=C_{m+n}^k.
  \]
  \item[б)] Докажите алгебраически, что
  \[
    (C_n^0)^2+(C_n^1)^2+\dots+(C_n^n)^2=C_{2n}^n.
  \]
\end{enumerate}

\item Докажите тождества (во втором пункте $n\geq2$):
\begin{enumerate}
  \item[а)] $C_n^1+2C_n^2+\dots+nC_n^n=n2^{n-1}$;
  \item[б)] $C_n^1-2C_n^2+\dots+(-1)^{n-1}nC_n^n=0$.
\end{enumerate}

\item Докажите, что булеву функцию $f(x_1,\ldots,x_n)$ можно выразить через $\mathbin{\&}$, $\vee$, $\supset$ тогда и только тогда, когда в её СКНФ отсутствует дизъюнкт $(\neg x_1\vee\cdots\vee\neg x_n)$.

\item Существует ли полная система, состоящая из одной булевой функции?

\item Сколько имеется четырёхзначных чисел, у которых каждая следующая цифра меньше предыдущей?

\item Дано число $53\underbrace{00\ldots00}_{100\text{ нулей}}35$. Требуется заменить некоторые два нуля на ненулевые цифры так, чтобы получилось число, делящееся на $495$. Сколькими способами это можно сделать?

\item Вычислите сумму $\displaystyle\sum_{k=0}^{n} 2^k C_n^k$.

\item Докажите, что набор простых дизъюнктов выполним (то есть можно так подобрать значения переменных, чтобы все эти дизъюнкты стали истинными) тогда и только тогда, когда из них при помощи правила резолюции нельзя получить пустой дизъюнкт.

\item Каждый житель острова людоедов принадлежит к одному из двух племён --- рыцарей, которые всегда говорят правду, или лжецов, которые всегда лгут. Однажды $2013$ островитян встали в круг и каждый заявил: «Оба моих соседа --- лжецы». Тут началась перебранка, и одного из островитян съели. После этого оставшиеся $2012$ островитян снова встали в круг, возможно, в другом порядке, и каждый заявил: «Оба моих соседа не из моего племени». Кого съели --- рыцаря или лжеца?

\end{problems}

\end{document}